% Libro 3 · anexos: guía rápida, plan de estudio y simulacro final
\makeatletter
\cleardoublepage\renewcommand{\zuka@mundoactual}{Guía rápida}
\phantomsection\addcontentsline{toc}{chapter}{Guía rápida}
{\sffamily\fontsize{26}{30}\selectfont\firalight Guía rápida\par}\vspace{2pt}
{\color{zlime}\rule{18mm}{1.6pt}}\par\vspace{6pt}

\begin{multicols}{2}\small
\kicker{Antes de leer (1)}
\begin{itemize}
  \item Título, fuente, primeras frases, último párrafo.
  \item Dos preguntas propias al texto.
  \item Mapa: una palabra de función por párrafo.
  \item EXCEPTO: tacha las que el texto dice.
\end{itemize}
\kicker{Qué texto es (2)}
\begin{itemize}
  \item Tipología: ¿para qué? (narrar, describir, explicar, argumentar, instruir).
  \item Género: cuento, crónica, noticia, ensayo, reseña, editorial, carta.
  \item Estructura: cronológica, problema-solución, causa-efecto, comparación.
\end{itemize}
\kicker{Lo que dice (3)}
\begin{itemize}
  \item Palabra ancla y lectura fina.
  \item Paráfrasis: conserva quién, cuántos, cuán seguro.
  \item Tres cajas: dice / niega / no dice.
  \item Un color por voz.
\end{itemize}
\columnbreak
\kicker{Palabras (4)}
\begin{itemize}
  \item Prueba del reemplazo.
  \item El contexto elige el sentido (literal o figurado).
  \item Frases hechas: nunca literales.
  \item Desarma: prefijo, raíz, sufijo.
\end{itemize}
\kicker{Hilos del texto (5, 6)}
\begin{itemize}
  \item Pronombre: concordancia y reemplazo.
  \item Elipsis: repón el verbo que falta.
  \item Conector: primero la relación.
  \item Función: un verbo (introduce, objeta, concluye…).
\end{itemize}
\kicker{Distractores}
\begin{itemize}
  \item Dice de más · Dice otra cosa · No lo dice · Dice lo contrario.
\end{itemize}
\end{multicols}

\clearpage\renewcommand{\zuka@mundoactual}{Plan de estudio}
\phantomsection\addcontentsline{toc}{chapter}{Plan de estudio con la app}
{\sffamily\fontsize{26}{30}\selectfont\firalight Plan de seis semanas\par}\vspace{2pt}
{\color{zlime}\rule{18mm}{1.6pt}}\par\vspace{6pt}
{\small Cada día: una lección del libro (lee el texto con lápiz, haz las preguntas comentadas tapando los comentarios) y
la misma lección en Zuka. Los sábados, la evaluación del mundo con reloj: 6 minutos de lectura y un minuto por pregunta.\par}
\vspace{6pt}
\begin{center}\small
\begin{tabularx}{\linewidth}{@{}>{\sffamily\bfseries\color{zolive}}lX>{\raggedleft\arraybackslash}p{24mm}@{}}
\toprule
\textbf{Semana} & \textbf{Lecciones} & \textbf{Evaluación}\\\midrule
1 & 1.1 Prelectura · 1.2 Idea central · 1.3 EXCEPTO & Mundo 1\\
2 & 2.1 Tipologías · 2.2 Géneros · 2.3 Estructura y testimonio & Mundo 2\\
3 & 3.1 Ubicar datos · 3.2 Paráfrasis · 3.3 Dice o no dice · 3.4 Voces & Mundo 3\\
4 & 4.1 Sinónimos · 4.2 Varios sentidos · 4.3 Expresiones · 4.4 Raíces & Mundo 4\\
5 & 5.1 Pronombres · 5.2 Elipsis · 6.1 Conectores · 6.2 Función & Mundos 5 y 6\\
6 & Simulacro final · repaso de los errores anotados & Simulacro\\
\bottomrule
\end{tabularx}
\end{center}

\clearpage\renewcommand{\zuka@mundoactual}{Simulacro final}
\setcounter{ej}{0}
\phantomsection\addcontentsline{toc}{chapter}{Simulacro final}
\immediate\write\zuka@sol{\noexpand\solmun{Simulacro final}}
\cabecera{SIMULACRO FINAL · LIBRO 3}{Leer lo que dice}{un texto · 10 preguntas\\16 minutos}\vspace{4pt}
\begin{texto}{El silencio de los pájaros}{Artículo de divulgación · texto original escrito para este libro}
\pa Quien se levanta temprano en una ciudad puede notar que los pájaros cantan antes del amanecer. No es casualidad: a esa
hora el aire está quieto, hay poco ruido y el sonido viaja más lejos. Para un ave que quiere marcar su territorio o atraer
pareja, es el mejor momento para hacerse oír.

\pa Sin embargo, el ruido de las ciudades ha cambiado esa costumbre. Diversos estudios han observado que algunas especies
urbanas cantan más agudo, más fuerte o a otras horas que sus parientes del campo. En otras palabras, se adaptan para no
quedar tapadas por el tráfico, del mismo modo que una persona levanta la voz en un restaurante lleno.

\pa Esa adaptación tiene un costo. Cantar más fuerte gasta más energía, y un canto modificado puede resultar menos
atractivo para las hembras, que en muchas especies eligen pareja según la calidad del canto. Además, no todas las especies
logran ajustarse: las que no lo consiguen tienden a desaparecer de los barrios más ruidosos.

\pa Durante los periodos en que el tráfico disminuyó de manera notable, investigadores de varios países registraron que
algunos pájaros volvieron a cantar más suave, como si de pronto ya no necesitaran gritar. Esos registros sugieren que los
cambios no son permanentes, sino una respuesta al ambiente.

\pa La conclusión es incómoda y esperanzadora a la vez. Incómoda, porque muestra que nuestro ruido afecta la vida de otros
seres de formas que no vemos. Esperanzadora, porque indica que, si bajamos el volumen, ellos pueden recuperar su voz.
\end{texto}
\begin{ejercicios}
\item ¿Cuál es la idea central del texto?
\opcionesl{Los pájaros cantan antes del amanecer.}{El ruido urbano modifica el canto de algunas aves, con costos, pero el cambio puede revertirse.}{Las hembras prefieren cantos agudos.}{El tráfico debe prohibirse.}
\sol{B. (1.2)}
\item Según el \S\,1, los pájaros cantan antes del amanecer porque
\opcionesl{tienen hambre.}{a esa hora el sonido viaja más lejos.}{las hembras duermen.}{el tráfico es mayor.}
\sol{B. (3.1)}
\item El texto es principalmente
\opcionesd{narrativo}{expositivo}{instructivo}{dialogal}
\sol{B. (2.1)}
\item La comparación con «una persona levanta la voz en un restaurante lleno» \pf{2} sirve para
\opcionesl{criticar a los restaurantes.}{hacer comprensible la adaptación de las aves.}{dar un dato científico.}{contradecir lo anterior.}
\sol{B. (6.2)}
\item «En otras palabras» \pf{2} introduce
\opcionesd{una reformulación}{una objeción}{una causa}{un ejemplo histórico}
\sol{A. (6.1)}
\item Según el \S\,3, todas son consecuencias o costos de la adaptación EXCEPTO
\opcionesl{mayor gasto de energía.}{un canto que puede ser menos atractivo.}{la desaparición de especies que no se adaptan.}{el aumento del número de pájaros en la ciudad.}
\sol{D. (1.3)}
\item En el \S\,3, «las que no lo consiguen» se refiere a
\opcionesd{las hembras}{las especies que no logran ajustarse}{las ciudades}{las parejas}
\sol{B. (5.1)}
\item ¿Cuál es una paráfrasis correcta de la idea del \S\,4?
\opcionesl{Los pájaros dejaron de cantar para siempre.}{Cuando hubo menos tráfico, algunos pájaros cantaron más suave, lo que sugiere que el cambio no es definitivo.}{Todos los pájaros cantan suave en el campo.}{Los investigadores obligaron a los pájaros a cantar.}
\sol{B. (3.2)}
\item El texto NO afirma que
\opcionesl{en muchas especies las hembras eligen pareja por el canto.}{algunas especies urbanas cantan más agudo.}{los pájaros del campo viven más años.}{el ruido afecta a otros seres de formas que no vemos.}
\sol{C. (3.3)}
\item En el \S\,5, la palabra «esperanzadora» se justifica porque
\opcionesl{los pájaros se están extinguiendo.}{bajar el ruido permitiría a las aves recuperar su canto.}{el ruido es incómodo.}{las ciudades crecen.}
\sol{B. (4.1)}
\end{ejercicios}
\makeatother
